% GENERATED FILE. Edit canonical lessons, then run npm run generate:books.
% canonical-source-hash: fnv1a64:2ca3fc40a4ab2028
% canonical-lessons: HI-C396-samskar, HI-C396-adars, HI-C396-udar, HI-C396-punya, HI-C396-ratna

\chapter{Values Passed On}
\label{ch:hi-values-passed-on}
\hlchaptermodality{396}

\begin{chapteropening}
\textbf{By the end of this chapter:} \emph{I can say what children will become if they are brought up well.}

Complete the last lesson of chapter 396: I can say what children will become if they are brought up well.
\end{chapteropening}

\section[\texorpdfstring{\hi{संस्कार} (saṁskār)}{संस्कार (saṁskār)}]{\hi{संस्कार} (saṁskār) — upbringing, values instilled}

\label{lesson:HI-C396-samskar}



\begin{quote}
Before the new one: say the Hindi for to fail, then the Hindi for in the least.
\end{quote}

\subsection*{You'll want to know: \hi{संस्कार}}
\textbf{\hi{संस्कार}} — \emph{saṁskār} — \textquotedblleft{}upbringing, values instilled\textquotedblright{}.

\begin{quote}
\hi{अगर} \hi{बच्चों} \hi{को} \hi{अच्छे} \hi{संस्कार} \hi{मिलेंगे} \hi{तो} \hi{वे} \hi{अच्छे} \hi{इंसान} \hi{बनेंगे।}
\end{quote}

\begin{quote}
\emph{agar baccoṁ ko acche saṁskār mileṅge to ve acche insān baneṅge.}
\end{quote}

\begin{quote}
If children are brought up well, they will become good people.
\end{quote}

Masculine.

\begin{cousinweb}
\textbf{\hi{संस्कार}} (\emph{saṁskār}) is a tatsama, borrowed from Sanskrit \emph{saṃskāra}, \textquotedblleft{}refining; an impression\textquotedblright{} (Monier-Williams).
\end{cousinweb}

\begin{grammarlens}[title={What you've built}]
One more word for saying what happens if.
\end{grammarlens}

\subsection*{Your turn}
\begin{itemize}
\raggedright
  \item \emph{Say it:} \emph{saṁskār}
  \item \emph{Say it:} \emph{saṁskār}, once more
  \item \emph{Say it:} \emph{agar baccoṁ ko acche saṁskār mileṅge to ve acche insān baneṅge.} — If children are brought up well, they will become good people.
  \item \emph{Recall:} say the Hindi for to fail, then the Hindi for in the least, then say \emph{saṁskār} again
\end{itemize}

\begin{tcolorbox}[breakable,colback=teal!4,colframe=teal!35!black,title={Before you move on}]
What does \hi{संस्कार} mean? (\textquotedblleft{}Upbringing, values instilled\textquotedblright{}.) Say it once more.
\end{tcolorbox}
\section[\texorpdfstring{\hi{आदर्श} (ādarś)}{आदर्श (ādarś)}]{\hi{आदर्श} (ādarś) — an ideal, a model}

\label{lesson:HI-C396-adars}



\begin{quote}
Before the new one: say the Hindi for in the least, then the Hindi for upbringing.
\end{quote}

\subsection*{You'll want to know: \hi{आदर्श}}
\textbf{\hi{आदर्श}} — \emph{ādarś} — \textquotedblleft{}an ideal, a model\textquotedblright{}.

\begin{quote}
\hi{अगर} \hi{तुम} \hi{सच} \hi{बोलोगे} \hi{तो} \hi{बच्चों} \hi{के} \hi{लिए} \hi{आदर्श} \hi{बनोगे।}
\end{quote}

\begin{quote}
\emph{agar tum sac bologe to baccoṁ ke liye ādarś banoge.}
\end{quote}

\begin{quote}
If you tell the truth, you'll be a model for the children.
\end{quote}

Masculine.

\begin{cousinweb}
\textbf{\hi{आदर्श}} (\emph{ādarś}) is a tatsama, borrowed from Sanskrit \emph{ādarśa}, \textquotedblleft{}a mirror\textquotedblright{} (Monier-Williams). Hindi uses it for a model to copy.
\end{cousinweb}

\begin{grammarlens}[title={What you've built}]
One more word for saying what happens if.
\end{grammarlens}

\subsection*{Your turn}
\begin{itemize}
\raggedright
  \item \emph{Say it:} \emph{ādarś}
  \item \emph{Say it:} \emph{ādarś}, once more
  \item \emph{Say it:} \emph{agar tum sac bologe to baccoṁ ke liye ādarś banoge.} — If you tell the truth, you'll be a model for the children.
  \item \emph{Recall:} say the Hindi for in the least, then the Hindi for upbringing, then say \emph{ādarś} again
\end{itemize}

\begin{tcolorbox}[breakable,colback=teal!4,colframe=teal!35!black,title={Before you move on}]
What does \hi{आदर्श} mean? (\textquotedblleft{}An ideal, a model\textquotedblright{}.) Say it once more.
\end{tcolorbox}
\section[\texorpdfstring{\hi{उदार} (udār)}{उदार (udār)}]{\hi{उदार} (udār) — generous, broad-minded}

\label{lesson:HI-C396-udar}



\begin{quote}
Before the new one: say the Hindi for upbringing, then the Hindi for an ideal.
\end{quote}

\subsection*{You'll want to know: \hi{उदार}}
\textbf{\hi{उदार}} — \emph{udār} — \textquotedblleft{}generous, broad-minded\textquotedblright{}.

\begin{quote}
\hi{अगर} \hi{वह} \hi{उदार} \hi{होता} \hi{तो} \hi{ग़रीबों} \hi{की} \hi{मदद} \hi{करता।}
\end{quote}

\begin{quote}
\emph{agar vah udār hotā to ġarīboṁ kī madad kartā.}
\end{quote}

\begin{quote}
If he were generous, he would help the poor.
\end{quote}

\begin{cousinweb}
\textbf{\hi{उदार}} (\emph{udār}) is a tatsama, borrowed from Sanskrit \emph{udāra}, \textquotedblleft{}great, noble\textquotedblright{} (Monier-Williams).
\end{cousinweb}

\begin{grammarlens}[title={What you've built}]
One more word for saying what happens if.
\end{grammarlens}

\subsection*{Your turn}
\begin{itemize}
\raggedright
  \item \emph{Say it:} \emph{udār}
  \item \emph{Say it:} \emph{udār}, once more
  \item \emph{Say it:} \emph{agar vah udār hotā to ġarīboṁ kī madad kartā.} — If he were generous, he would help the poor.
  \item \emph{Recall:} say the Hindi for upbringing, then the Hindi for an ideal, then say \emph{udār} again
\end{itemize}

\begin{tcolorbox}[breakable,colback=teal!4,colframe=teal!35!black,title={Before you move on}]
What does \hi{उदार} mean? (\textquotedblleft{}Generous, broad-minded\textquotedblright{}.) Say it once more.
\end{tcolorbox}
\section[\texorpdfstring{\hi{पुण्य} (puṇya)}{पुण्य (puṇya)}]{\hi{पुण्य} (puṇya) — a good deed, merit}

\label{lesson:HI-C396-punya}



\begin{quote}
Before the new one: say the Hindi for an ideal, then the Hindi for generous.
\end{quote}

\subsection*{You'll want to know: \hi{पुण्य}}
\textbf{\hi{पुण्य}} — \emph{puṇya} — \textquotedblleft{}a good deed, merit\textquotedblright{}.

\begin{quote}
\hi{अगर} \hi{तुम} \hi{भूखे} \hi{को} \hi{खाना} \hi{दोगे} \hi{तो} \hi{पुण्य} \hi{मिलेगा।}
\end{quote}

\begin{quote}
\emph{agar tum bhūkhe ko khānā doge to puṇya milegā.}
\end{quote}

\begin{quote}
If you give food to the hungry, you will earn merit.
\end{quote}

Masculine.

\begin{cousinweb}
\textbf{\hi{पुण्य}} (\emph{puṇya}) is a tatsama, borrowed from Sanskrit \emph{puṇya}, \textquotedblleft{}good, auspicious; merit\textquotedblright{} (Monier-Williams).
\end{cousinweb}

\begin{grammarlens}[title={What you've built}]
One more word for saying what happens if.
\end{grammarlens}

\subsection*{Your turn}
\begin{itemize}
\raggedright
  \item \emph{Say it:} \emph{puṇya}
  \item \emph{Say it:} \emph{puṇya}, once more
  \item \emph{Say it:} \emph{agar tum bhūkhe ko khānā doge to puṇya milegā.} — If you give food to the hungry, you will earn merit.
  \item \emph{Recall:} say the Hindi for an ideal, then the Hindi for generous, then say \emph{puṇya} again
\end{itemize}

\begin{tcolorbox}[breakable,colback=teal!4,colframe=teal!35!black,title={Before you move on}]
What does \hi{पुण्य} mean? (\textquotedblleft{}A good deed, merit\textquotedblright{}.) Say it once more.
\end{tcolorbox}
\section[\texorpdfstring{\hi{रटना} (raṭnā)}{रटना (raṭnā)}]{\hi{रटना} (raṭnā) — to learn by rote, to cram}

\label{lesson:HI-C396-ratna}



\begin{quote}
Before the new one: say the Hindi for generous, then the Hindi for a good deed.
\end{quote}

\subsection*{You'll want to know: \hi{रटना}}
\textbf{\hi{रटना}} — \emph{raṭnā} — \textquotedblleft{}to learn by rote, to cram\textquotedblright{}.

\begin{quote}
\hi{अगर} \hi{तुम} \hi{केवल} \hi{रटोगे} \hi{तो} \hi{कुछ} \hi{नहीं} \hi{समझोगे।}
\end{quote}

\begin{quote}
\emph{agar tum keval raṭoge to kuch nahīṃ samjhoge.}
\end{quote}

\begin{quote}
If you only learn by rote, you won't understand anything.
\end{quote}

Past, with \hi{ने} (he or she …): \hi{उसने} \hi{पूरा} \hi{पाठ} \hi{रटा}, \emph{usne pūrā pāṭh raṭā}.

\begin{cousinweb}
\textbf{\hi{रटना}} (\emph{raṭnā}) is a tadbhava, inherited from Sanskrit \emph{raṭati}, \textquotedblleft{}cries, yells\textquotedblright{} (CDIAL 10590). Hindi uses the sense of saying over and over.
\end{cousinweb}

\begin{grammarlens}[title={What you've built}]
One more word for saying what happens if.
\end{grammarlens}

\subsection*{Your turn}
\begin{itemize}
\raggedright
  \item \emph{Say it:} \emph{raṭnā}
  \item \emph{Say it:} \emph{raṭnā}, once more
  \item \emph{Say it:} \emph{agar tum keval raṭoge to kuch nahīṃ samjhoge.} — If you only learn by rote, you won't understand anything.
  \item \emph{Recall:} say the Hindi for generous, then the Hindi for a good deed, then say \emph{raṭnā} again
\end{itemize}

\begin{tcolorbox}[breakable,colback=teal!4,colframe=teal!35!black,title={Before you move on}]
What does \hi{रटना} mean? (\textquotedblleft{}To learn by rote, to cram\textquotedblright{}.) Say it once more.
\end{tcolorbox}
